\documentclass[12pt,a4paper]{article}
\usepackage[utf-8]{inputenc}
\usepackage[french]{babel}
\usepackage[margin=1in]{geometry}
\usepackage{xcolor}
\usepackage{listings}
\usepackage{fancyhdr}
\usepackage{hyperref}
\usepackage{graphicx}
\usepackage{amsmath}
\usepackage{fancybox}
\usepackage{tikz}
\usetikzlibrary{shapes,arrows,positioning}

% Configuration des couleurs
\definecolor{codebackground}{gray}{0.95}
\definecolor{codekeyword}{rgb}{0.13,0.29,0.53}
\definecolor{codestring}{rgb}{0.63,0.13,0.94}
\definecolor{codecomment}{rgb}{0.13,0.54,0.13}

% Configuration du style de code
\lstset{
    backgroundcolor=\color{codebackground},
    basicstyle=\ttfamily\small,
    breaklines=true,
    commentstyle=\color{codecomment},
    keywordstyle=\color{codekeyword},
    stringstyle=\color{codestring},
    showstringspaces=false,
    tabsize=4,
    frame=single,
    rulecolor=\color{black},
    captionpos=b
}

% En-tête et pied de page
\pagestyle{fancy}
\lhead{\textbf{EduPredict}}
\rhead{\textit{Guide Complet}}
\cfoot{\thepage}

\title{\textbf{📚 EduPredict} \\ \Large Guide Complet : Installation, Configuration et Liaison API-Laravel}
\author{Documentation Technique}
\date{\today}

\begin{document}

\maketitle

\tableofcontents
\newpage

% ==================== SECTION 1 ====================
\section{Vue d'ensemble du projet}

\subsection{Architecture générale}

Le projet \textbf{EduPredict} est une application web de prédiction de réussite scolaire composée de deux composants principaux :

\begin{enumerate}
    \item \textbf{Backend API (Python FastAPI)} : API de prédiction basée sur un modèle de machine learning
    \item \textbf{Frontend Web (Laravel + Blade)} : Interface web avec authentification et tableau de bord
\end{enumerate}

\subsection{Flux de communication}

\begin{figure}[h]
\centering
\begin{tikzpicture}[node distance=2cm]
    % Nœuds
    \node (user) [rectangle, draw, fill=blue!20] {Utilisateur};
    \node (laravel) [rectangle, draw, fill=green!20, right of=user, xshift=2cm] {Laravel};
    \node (fastapi) [rectangle, draw, fill=orange!20, right of=laravel, xshift=2cm] {FastAPI};
    \node (model) [rectangle, draw, fill=red!20, right of=fastapi, xshift=2cm] {ML Model};
    
    % Flèches
    \draw[->] (user) -- (laravel) node[midway, above] {Formulaire};
    \draw[->] (laravel) -- (fastapi) node[midway, above] {HTTP POST};
    \draw[->] (fastapi) -- (model) node[midway, above] {Prédiction};
    \draw[->] (model) -- (fastapi) node[midway, below] {Résultat};
    \draw[->] (fastapi) -- (laravel) node[midway, below] {JSON};
    \draw[->] (laravel) -- (user) node[midway, below] {Affichage};
\end{tikzpicture}
\caption{Architecture générale du système}
\end{figure}

\newpage

% ==================== SECTION 2 ====================
\section{Installation et Configuration}

\subsection{Prérequis système}

\begin{itemize}
    \item \textbf{Python 3.9+} : Pour l'API FastAPI
    \item \textbf{PHP 8.1+} : Pour Laravel
    \item \textbf{Composer} : Gestionnaire de paquets PHP
    \item \textbf{Node.js} : Pour les assets frontend
    \item \textbf{SQLite/MySQL} : Base de données
\end{itemize}

\subsection{Étape 1 : Installation de Laravel}

\subsubsection{1.1 Cloner le projet}

\begin{lstlisting}[language=bash]
cd c:\Users\user\Desktop\PROJETTTTTTTTT\wafamay\classrrom
cd edupredict
\end{lstlisting}

\subsubsection{1.2 Installer les dépendances PHP}

\begin{lstlisting}[language=bash]
# Installation des paquets Composer
composer install

# Installation des paquets npm
npm install
\end{lstlisting}

\subsubsection{1.3 Créer le fichier .env}

\begin{lstlisting}[language=bash]
# Créer une copie du fichier d'exemple
copy .env.example .env

# Générer la clé d'application
php artisan key:generate
\end{lstlisting}

\subsubsection{1.4 Configuration du fichier .env}

Modifier le fichier \texttt{.env} avec les valeurs appropriées :

\begin{lstlisting}[style=,basicstyle=\ttfamily\footnotesize]
APP_NAME=EduPredict
APP_ENV=local
APP_DEBUG=true
APP_KEY=base64:...

# Configuration de la base de données
DB_CONNECTION=sqlite
# DB_HOST=127.0.0.1
# DB_PORT=3306
# DB_DATABASE=edupredict
# DB_USERNAME=root
# DB_PASSWORD=

# Configuration de l'API de prédiction
PREDICT_API_URL=http://127.0.0.1:8001/predict
PREDICT_API_START_COMMAND=python -m uvicorn api:app --host 127.0.0.1 --port 8001
\end{lstlisting}

\subsubsection{1.5 Configuration de la base de données}

\begin{lstlisting}[language=bash]
# Exécuter les migrations
php artisan migrate

# (Optionnel) Remplir la BD avec des données de test
php artisan db:seed
\end{lstlisting}

\subsubsection{1.6 Compiler les assets}

\begin{lstlisting}[language=bash]
# Mode développement (avec hot reload)
npm run dev

# Mode production
npm run build
\end{lstlisting}

\subsection{Étape 2 : Installation et configuration de l'API Python}

\subsubsection{2.1 Accéder au dossier de l'API}

\begin{lstlisting}[language=bash]
cd ..
cd student_project\student_project
\end{lstlisting}

\subsubsection{2.2 Créer un environnement virtuel Python}

\begin{lstlisting}[language=bash]
# Créer l'environnement virtuel
python -m venv venv

# Activer l'environnement
# Sous Windows (PowerShell)
.\venv\Scripts\Activate.ps1

# Sous Windows (CMD)
venv\Scripts\activate

# Sous Linux/macOS
source venv/bin/activate
\end{lstlisting}

\subsubsection{2.3 Installer les dépendances Python}

\begin{lstlisting}[language=bash]
# Mettre à jour pip
pip install --upgrade pip

# Installer les dépendances
pip install fastapi uvicorn joblib pandas scikit-learn
\end{lstlisting}

\subsubsection{2.4 Vérifier les fichiers du modèle}

\begin{lstlisting}[language=bash]
# Les fichiers suivants doivent exister dans le dossier "model/"
ls model/
# Output:
# model.pkl
# scaler.pkl
\end{lstlisting}

\newpage

% ==================== SECTION 3 ====================
\section{Démarrage et Exécution}

\subsection{Démarrage de l'API FastAPI}

\subsubsection{3.1 Premier démarrage}

\begin{lstlisting}[language=bash]
# Depuis le répertoire student_project/student_project
# S'assurer que l'environnement virtuel est activé

python -m uvicorn api:app --host 127.0.0.1 --port 8001 --reload
\end{lstlisting}

\textbf{Résultat attendu :}
\begin{lstlisting}[style=,basicstyle=\ttfamily\footnotesize]
INFO:     Uvicorn running on http://127.0.0.1:8001
INFO:     Application startup complete
\end{lstlisting}

\subsubsection{3.2 Vérifier l'API}

Ouvrir un navigateur et visiter : \url{http://127.0.0.1:8001}

La réponse doit être :
\begin{lstlisting}[language=json]
{
  "message": "Student API working"
}
\end{lstlisting}

\subsection{Démarrage de Laravel}

\subsubsection{3.3 Lancer le serveur Laravel}

\begin{lstlisting}[language=bash]
# Depuis le répertoire edupredict
php artisan serve
\end{lstlisting}

\textbf{Résultat attendu :}
\begin{lstlisting}[style=,basicstyle=\ttfamily\footnotesize]
INFO Server running on [http://127.0.0.1:8000]
\end{lstlisting}

\subsubsection{3.4 Accéder à l'application}

Ouvrir un navigateur et visiter : \url{http://127.0.0.1:8000}

\subsection{Commandes complètes pour démarrage rapide}

\textbf{Terminal 1 - API FastAPI :}
\begin{lstlisting}[language=bash]
cd c:\Users\user\Desktop\PROJETTTTTTTTT\wafamay\classrrom\student_project\student_project
.\venv\Scripts\Activate.ps1
python -m uvicorn api:app --host 127.0.0.1 --port 8001 --reload
\end{lstlisting}

\textbf{Terminal 2 - Serveur de développement Node :}
\begin{lstlisting}[language=bash]
cd c:\Users\user\Desktop\PROJETTTTTTTPT\wafamay\classrrom\edupredict
npm run dev
\end{lstlisting}

\textbf{Terminal 3 - Application Laravel :}
\begin{lstlisting}[language=bash]
cd c:\Users\user\Desktop\PROJETTTTTTTPT\wafamay\classrrom\edupredict
php artisan serve
\end{lstlisting}

\newpage

% ==================== SECTION 4 ====================
\section{Liaison API-Laravel : Explication Détaillée}

\subsection{4.1 Configuration de l'API dans Laravel}

La configuration de l'API est centralisée dans le fichier \texttt{config/services.php} :

\begin{lstlisting}[language=php]
'predict_api' => [
    'url' => env('PREDICT_API_URL', 'http://127.0.0.1:8001/predict'),
    'start_command' => env('PREDICT_API_START_COMMAND', 
        'python -m uvicorn api:app --host 127.0.0.1 --port 8001'),
],
\end{lstlisting}

\subsection{4.2 Flow de la requête utilisateur}

\subsubsection{Étape 1 : Soumission du formulaire}

L'utilisateur complète le formulaire de prédiction :

\begin{lstlisting}[language=html]
<form method="POST" action="{{ route('student.predict') }}">
    @csrf
    <input name="name" type="text" required>
    <input name="age" type="number" min="15" max="25" required>
    <input name="study_time" type="number" min="1" max="4" required>
    <input name="failures" type="number" min="0" max="3" required>
    <input name="absence" type="number" min="0" required>
    <input name="filiere" type="text">
    <button type="submit">Prédire</button>
</form>
\end{lstlisting}

\subsubsection{Étape 2 : Traitement dans le contrôleur Laravel}

La route dirige vers \texttt{StudentController@predict} :

\begin{lstlisting}[language=php]
// Routes (routes/web.php)
Route::post('/student/predict', 
    [StudentController::class, 'predict']
)->name('student.predict');
\end{lstlisting}

\subsubsection{Étape 3 : Validation des données}

Le contrôleur valide les données reçues :

\begin{lstlisting}[language=php]
// app/Http/Controllers/StudentController.php
public function predict(Request $request): RedirectResponse
{
    $validated = $request->validate([
        'name'       => ['required', 'string', 'max:255'],
        'age'        => ['required', 'integer', 'min:15', 'max:25'],
        'study_time' => ['required', 'integer', 'min:1', 'max:4'],
        'failures'   => ['required', 'integer', 'min:0', 'max:3'],
        'absence'    => ['required', 'integer', 'min:0'],
        'filiere'    => ['nullable', 'string', 'max:255'],
    ]);
    // ...
}
\end{lstlisting}

\subsubsection{Étape 4 : Appel HTTP à l'API FastAPI}

Laravel effectue une requête POST vers l'API :

\begin{lstlisting}[language=php]
try {
    $response = Http::timeout(10)
        ->acceptJson()
        ->post(config('services.predict_api.url'), [
            'age'      => $validated['age'],
            'studytime' => $validated['study_time'],
            'failures' => $validated['failures'],
            'absences' => $validated['absence'],
        ])
        ->throw()
        ->json();
} catch (ConnectionException $exception) {
    return back()
        ->withInput()
        ->withErrors([
            'api' => "Impossible de joindre l'API..."
        ]);
}
\end{lstlisting}

\textbf{Détails de la requête :}
\begin{itemize}
    \item \textbf{Méthode} : POST
    \item \textbf{URL} : \texttt{http://127.0.0.1:8001/predict}
    \item \textbf{Timeout} : 10 secondes
    \item \textbf{Corps} : JSON avec les données du formulaire
    \item \textbf{Headers} : \texttt{Accept: application/json}
\end{itemize}

\subsubsection{Étape 5 : Réponse de l'API FastAPI}

L'API FastAPI reçoit la requête et traite les données :

\begin{lstlisting}[language=python]
# student_project/student_project/api.py
@app.post("/predict")
def predict(data: dict):
    try:
        # 1. Convertir les données en DataFrame
        df = pd.DataFrame([data])
        
        # 2. Normaliser les données
        scaled_features = scaler.transform(df)
        
        # 3. Faire la prédiction (0 ou 1)
        prediction = int(model.predict(scaled_features)[0])
        
        # 4. Calculer la probabilité
        probability = float(
            model.predict_proba(scaled_features)[0][1]
        )
        
        # 5. Retourner le résultat en JSON
        return {
            "prediction": prediction,
            "probability_success": probability,
        }
    except Exception as exc:
        raise HTTPException(
            status_code=500, 
            detail=str(exc)
        )
\end{lstlisting}

\subsubsection{Étape 6 : Traitement de la réponse}

Laravel traite la réponse de l'API :

\begin{lstlisting}[language=php]
$prediction = (int) ($response['prediction'] ?? 0);
$probability = isset($response['probability_success'])
    ? round(((float) $response['probability_success']) * 100, 2)
    : null;
$result = $prediction === 1 ? 'Reussi' : 'Echoue';
\end{lstlisting}

\subsubsection{Étape 7 : Enregistrement en base de données}

Les résultats sont stockés :

\begin{lstlisting}[language=php]
Student::create([
    'user_id'    => $request->user()->id,
    'name'       => $validated['name'],
    'age'        => $validated['age'],
    'study_time' => $validated['study_time'],
    'failures'   => $validated['failures'],
    'absence'    => $validated['absence'],
    'filiere'    => $validated['filiere'] ?? '',
    'probability' => $probability,
    'result'     => $result,
]);
\end{lstlisting}

\subsubsection{Étape 8 : Retour à l'utilisateur}

L'application redirige avec les résultats :

\begin{lstlisting}[language=php]
return redirect()
    ->route('student.dashboard')
    ->with([
        'result'      => $result,
        'probability' => $probability,
    ]);
\end{lstlisting}

\newpage

\subsection{4.3 Diagramme de flux complet}

\begin{figure}[h]
\centering
\begin{tikzpicture}[node distance=2.5cm, thick]
    % Styles
    \tikzstyle{usernode} = [rectangle, draw, fill=blue!20, text width=2cm, text centered]
    \tikzstyle{laravelnode} = [rectangle, draw, fill=green!20, text width=2cm, text centered]
    \tikzstyle{apinode} = [rectangle, draw, fill=orange!20, text width=2cm, text centered]
    \tikzstyle{dbnode} = [cylinder, draw, fill=purple!20, text width=1.5cm, text centered, aspect=0.4]
    
    % Nœuds
    \node (user) [usernode] {Utilisateur \\ remplit formulaire};
    \node (valid) [laravelnode, below of=user] {Laravel valide \\ les données};
    \node (apicall) [laravelnode, below of=valid] {POST à l'API \\ /predict};
    \node (apireceive) [apinode, below of=apicall] {FastAPI reçoit \\ les données};
    \node (process) [apinode, below of=apireceive] {Normalisation \\ + Prédiction};
    \node (response) [apinode, below of=process] {Retour JSON \\ (résultat + prob)};
    \node (store) [dbnode, right of=valid, xshift=3.5cm] {Base de \\ données};
    \node (display) [laravelnode, below of=response] {Laravel affiche \\ les résultats};
    
    % Flèches
    \draw[->] (user) -- (valid);
    \draw[->] (valid) -- (apicall);
    \draw[->] (apicall) -- (apireceive);
    \draw[->] (apireceive) -- (process);
    \draw[->] (process) -- (response);
    \draw[->] (response) -- (display);
    \draw[->] (display) -- (store);
\end{tikzpicture}
\caption{Flux complet de la prédiction}
\end{figure}

\subsection{4.4 Gestion des erreurs}

\subsubsection{ConnectionException - API non accessible}

\begin{lstlisting}[language=php]
catch (ConnectionException $exception) {
    return back()
        ->withInput()
        ->withErrors([
            'api' => "Impossible de joindre l'API de prediction..."
        ]);
}
\end{lstlisting}

\textbf{Solutions :}
\begin{enumerate}
    \item Vérifier que l'API FastAPI est en cours d'exécution
    \item Vérifier l'URL dans \texttt{.env}
    \item Vérifier la connexion réseau
\end{enumerate}

\subsubsection{RequestException - API retourne une erreur}

\begin{lstlisting}[language=php]
catch (RequestException $exception) {
    return back()
        ->withInput()
        ->withErrors([
            'api' => "L'API de prediction a repondu avec une erreur..."
        ]);
}
\end{lstlisting}

\textbf{Solutions :}
\begin{enumerate}
    \item Vérifier les logs de l'API
    \item Vérifier les données envoyées
    \item Vérifier que le modèle ML est bien chargé
\end{enumerate}

\newpage

% ==================== SECTION 5 ====================
\section{Fichiers et Structure}

\subsection{Structure Laravel}

\begin{lstlisting}[style=,basicstyle=\ttfamily\footnotesize]
edupredict/
├── app/
│   ├── Http/
│   │   ├── Controllers/
│   │   │   ├── StudentController.php    # Gère les prédictions
│   │   │   ├── AdminController.php      # Gère le dashboard admin
│   │   │   └── AuthController.php       # Authentification
│   │   └── Middleware/
│   ├── Models/
│   │   ├── Student.php                  # Modèle des prédictions
│   │   └── User.php                     # Modèle utilisateurs
│   └── Providers/
├── config/
│   ├── app.php
│   ├── auth.php
│   ├── database.php
│   └── services.php                     # CONFIG API (important!)
├── database/
│   ├── migrations/                      # Schéma BD
│   └── seeders/                         # Données de test
├── resources/
│   ├── views/
│   │   ├── welcome.blade.php            # Page accueil
│   │   ├── admin/
│   │   ├── auth/                        # Pages login/register
│   │   └── student/                     # Pages student
│   ├── css/
│   └── js/
├── routes/
│   └── web.php                          # Routes (important!)
└── .env                                 # Variables d'environnement
\end{lstlisting}

\subsection{Structure API FastAPI}

\begin{lstlisting}[style=,basicstyle=\ttfamily\footnotesize]
student_project/student_project/
├── api.py                               # API FastAPI (important!)
├── app.py                               # App Flask (optionnel)
├── train.py                             # Script d'entraînement
├── model/
│   ├── model.pkl                        # Modèle ML
│   └── scaler.pkl                       # Normaliseur
├── data/                                # Données d'entraînement
└── venv/                                # Environnement Python
\end{lstlisting}

\newpage

% ==================== SECTION 6 ====================
\section{Commandes Utiles}

\subsection{Commandes Laravel}

\begin{table}[h]
\centering
\small
\begin{tabular}{|p{4cm}|p{8cm}|}
\hline
\textbf{Commande} & \textbf{Description} \\
\hline
\texttt{php artisan serve} & Démarrer le serveur dev \\
\hline
\texttt{php artisan migrate} & Exécuter les migrations \\
\hline
\texttt{php artisan migrate:refresh} & Réinitialiser la BD \\
\hline
\texttt{php artisan db:seed} & Remplir avec données de test \\
\hline
\texttt{php artisan tinker} & Console interactive \\
\hline
\texttt{php artisan route:list} & Voir toutes les routes \\
\hline
\texttt{npm run dev} & Compiler assets (dev) \\
\hline
\texttt{npm run build} & Compiler assets (prod) \\
\hline
\end{tabular}
\end{table}

\subsection{Commandes Python/API}

\begin{table}[h]
\centering
\small
\begin{tabular}{|p{4cm}|p{8cm}|}
\hline
\textbf{Commande} & \textbf{Description} \\
\hline
\texttt{python -m venv venv} & Créer environnement virtuel \\
\hline
\texttt{.\textbackslash venv\textbackslash Scripts\textbackslash Activate.ps1} & Activer venv (Windows PS) \\
\hline
\texttt{venv\textbackslash Scripts\textbackslash activate} & Activer venv (Windows CMD) \\
\hline
\texttt{pip install -r requirements.txt} & Installer dépendances \\
\hline
\texttt{python -m uvicorn api:app --reload} & Lancer API en dev \\
\hline
\texttt{python train.py} & Entraîner le modèle \\
\hline
\end{tabular}
\end{table}

\newpage

% ==================== SECTION 7 ====================
\section{Résumé des Étapes}

\subsection{Pour la première utilisation}

\begin{enumerate}
    \item \textbf{Cloner le projet}
    \item \textbf{Installer Laravel}
        \begin{itemize}
            \item \texttt{composer install}
            \item \texttt{npm install}
            \item \texttt{cp .env.example .env}
            \item \texttt{php artisan key:generate}
            \item \texttt{php artisan migrate}
        \end{itemize}
    \item \textbf{Installer l'API}
        \begin{itemize}
            \item Créer venv : \texttt{python -m venv venv}
            \item Activer : \texttt{.\textbackslash venv\textbackslash Scripts\textbackslash Activate.ps1}
            \item Installer : \texttt{pip install fastapi uvicorn joblib pandas scikit-learn}
        \end{itemize}
    \item \textbf{Lancer les services (3 terminaux)}
        \begin{itemize}
            \item Terminal 1 : API Python \texttt{python -m uvicorn api:app --reload}
            \item Terminal 2 : \texttt{npm run dev}
            \item Terminal 3 : \texttt{php artisan serve}
        \end{itemize}
    \item \textbf{Accéder à} \url{http://127.0.0.1:8000}
\end{enumerate}

\subsection{Fichiers clés à connaître}

\begin{itemize}
    \item \textbf{config/services.php} : Configuration API
    \item \textbf{routes/web.php} : Routes Laravel
    \item \textbf{app/Http/Controllers/StudentController.php} : Logique prédiction
    \item \textbf{student\_project/student\_project/api.py} : Endpoint API
    \item \textbf{.env} : Variables d'environnement
\end{itemize}

\newpage

% ==================== SECTION 8 ====================
\section{Dépannage}

\subsection{Problèmes courants}

\begin{enumerate}
    \item \textbf{Erreur : "Impossible de joindre l'API"}
    \begin{itemize}
        \item ✓ Vérifier que \texttt{python -m uvicorn api:app} est en cours d'exécution
        \item ✓ Vérifier que le port 8001 est libre
        \item ✓ Vérifier \texttt{PREDICT\_API\_URL} dans \texttt{.env}
    \end{itemize}
    
    \item \textbf{Erreur : "Model file not found"}
    \begin{itemize}
        \item ✓ Vérifier que \texttt{model/model.pkl} existe
        \item ✓ Vérifier que \texttt{model/scaler.pkl} existe
        \item ✓ Lancer \texttt{python train.py} pour générer les fichiers
    \end{itemize}
    
    \item \textbf{Erreur : "Connection refused"}
    \begin{itemize}
        \item ✓ Vérifier le firewall Windows
        \item ✓ Essayer avec localhost au lieu de 127.0.0.1
    \end{itemize}
    
    \item \textbf{Erreur de base de données}
    \begin{itemize}
        \item ✓ Lancer \texttt{php artisan migrate:refresh}
        \item ✓ Vérifier \texttt{DB\_CONNECTION} dans \texttt{.env}
    \end{itemize}
\end{enumerate}

\section{Conclusion}

Le projet EduPredict intègre une \textbf{API Python} et une \textbf{application web Laravel} à travers des appels HTTP. La communication se fait de manière synchrone : Laravel envoie les données utilisateur à l'API, attend la prédiction, puis stocke le résultat en base de données.

Cette architecture permet une séparation des responsabilités :
\begin{itemize}
    \item \textbf{API} : Gère uniquement la prédiction ML
    \item \textbf{Laravel} : Gère l'interface, l'authentification et la persistance
\end{itemize}

\end{document}
